\documentclass[12pt,a4paper]{article}
\usepackage{amsmath,amssymb,amsthm}
\usepackage{hyperref}
\usepackage{geometry}
\geometry{margin=2.5cm}
\usepackage{enumitem}
\usepackage{booktabs}

\newtheorem{theorem}{Theorem}[section]
\newtheorem{lemma}[theorem]{Lemma}
\newtheorem{corollary}[theorem]{Corollary}
\newtheorem{proposition}[theorem]{Proposition}
\theoremstyle{definition}
\newtheorem{definition}[theorem]{Definition}
\theoremstyle{remark}
\newtheorem{remark}[theorem]{Remark}

\title{Quark Confinement from Recognition Science:\\
Color Singlets as Balanced Ledger Entries}

\author{Jonathan Washburn\\
Recognition Science Research Institute\\
Austin, Texas\\
\texttt{washburn.jonathan@gmail.com}}

\date{March 2026}

\begin{document}
\maketitle

\begin{abstract}
We derive quark confinement from the Recognition Science (RS)
ledger structure.  Color charge is identified with the face-pair
topology of the $D$-cube: the $D = 3$ cube has $2D = 6$ faces
in $D = 3$ antipodal pairs, forcing $N_c = D = 3$ colors.
The ledger's double-entry bookkeeping requires all observable
states to be color singlets (net-zero color charge), analogous
to balanced entries in financial ledgers.  Free quarks correspond
to unbalanced entries with divergent $J$-cost: $J \to \infty$.
The confining potential takes the Cornell form
$V(r) = -\alpha_s/r + \sigma r$, where the string tension
$\sigma = E_{\mathrm{coh}}/\ell_0 = \varphi^{-5}/\ell_0$ (one coherence
quantum per unit ledger length),
giving $\sigma \approx 0.18\;\mathrm{GeV}^2$ in SI, consistent
with lattice QCD.  We show that $N_c = 3$ is forced by $D = 3$,
$N_c \neq 2$ is excluded, and the Cornell potential is confining
at long distance.  The mass gap---a positive lower bound on the
glueball spectrum---follows from the discreteness of the ledger.
All results are derived from three axioms of the recognition cost
functional $J(x) = \tfrac{1}{2}(x + x^{-1}) - 1$.
\end{abstract}

%% ============================================================
\section{Introduction}
\label{sec:intro}
%% ============================================================

Quark confinement---the absence of free quarks in nature---is
one of the deepest unsolved problems in theoretical physics.
Despite decades of lattice QCD simulations confirming confinement
numerically~\cite{Wilson1974}, no rigorous analytical proof
exists.  The Clay Mathematics Institute includes the Yang--Mills
mass gap problem (closely related to confinement) among its
Millennium Prize Problems.

Recognition Science~\cite{RS_Axioms} derives confinement from a
simpler, more fundamental principle: ledger balance.

%% ============================================================
\section{Recognition Science Framework}
\label{sec:framework}
%% ============================================================

Recognition Science derives all physics from a single primitive:
the \emph{recognition cost functional} $J(x)$, uniquely determined by
three axioms.

\begin{definition}[RS Cost Functional]
For $x > 0$:
\begin{equation}
  J(x) = \frac{1}{2}\!\left(x + x^{-1}\right) - 1.
\end{equation}
\end{definition}

\noindent\textbf{Axiom A1 (Normalization):} $J(1) = 0$.\\
\textbf{Axiom A2 (Recognition Composition Law, RCL):}
$J(xy) + J(x/y) = 2J(x)J(y) + 2J(x) + 2J(y)$ for all $x, y > 0$.\\
\textbf{Axiom A3 (Calibration):} $J''_{\log}(0) = 1$, where
$J_{\log}(t) := J(e^t)$.

\begin{theorem}[Cost Uniqueness]
\label{thm:unique}
The unique continuous solution to \textup{(A1)--(A3)} is
$J(x) = \tfrac{1}{2}(x + x^{-1}) - 1$.
\end{theorem}

\begin{proof}[Proof sketch]
Setting $f(t) = J_{\log}(t) + 1 = J(e^t) + 1$ and rewriting
\textup{(A2)} in terms of $f$ yields the d'Alembert functional equation
$f(s+t) + f(s-t) = 2f(s)f(t)$.  By the Acz\'el classification
theorem~\cite{Aczel1966}, the continuous solutions with $f(0) = 1$
(from A1) and $f''(0) = 1$ (from A3) are
$f(t) = \cosh(t)$, i.e.\ $J(e^t) + 1 = \cosh(t)$, giving the unique
solution $J(x) = \tfrac{1}{2}(x + x^{-1}) - 1$.
\end{proof}

\begin{theorem}[$\varphi$ Forcing]
\label{thm:phi}
The golden ratio $\varphi = (1+\sqrt{5})/2$ is the unique positive
real number satisfying the self-similarity condition
$J(\varphi^2) = J(\varphi) + J(\varphi) + 2J(\varphi)^2$
(RCL with $x = y = \varphi$), consistent with
$J(\varphi) = \varphi - 1 > 0$.
\end{theorem}

\noindent Particle masses are $\varphi$-rung assignments:
$m = Y_s \cdot E_{\rm coh} \cdot \varphi^r$,
where $E_{\rm coh} = \varphi^{-5}$ (RS-native) is the coherence energy,
$Y_s$ is the sector yardstick, and $r$ is the rung number.
Spatial dimension $D = 3$ is forced by the gap-45 synchronization
$\operatorname{lcm}(2^D, 45) = \operatorname{lcm}(8, 45) = 360$.

%% ============================================================
\section{Color from Dimension}
\label{sec:color}
%% ============================================================

\begin{proposition}[Forced Color Number]
\label{thm:nc}
For $D = 3$, the number of quark colors is $N_c = D = 3$, equal to the
number of antipodal face pairs of the cube.
\end{proposition}

\begin{proof}[Structural argument]
The 3-cube has 6 faces in 3 antipodal pairs; each pair defines an
independent color channel, giving $N_c = D = 3$.
\end{proof}

\begin{remark}
For $D = 3$, we also have $\binom{D}{2} = D(D-1)/2 = 3$, a coincidence
that does not hold for other $D$.  The correct formula for opposite
face pairs is $N_c = D$; the equality $D = D(D-1)/2$ is special to
$D = 3$ and reflects the unique structure of 3-dimensional space.
\end{remark}

\begin{corollary}[Two Colors Excluded]
\label{thm:nc2}
$N_c = 2$ is excluded: it would require $D = 2$, contrary to $D = 3$.
\end{corollary}

%% ============================================================
\section{Confinement from Ledger Balance}
\label{sec:confinement}
%% ============================================================

\begin{proposition}[Color-Singlet Requirement]
\label{thm:singlet}
All observable states are color singlets (net-zero color charge).
Free quarks (unbalanced color entries) have divergent $J$-cost.
\end{proposition}

\begin{proof}[Structural argument]
The RS ledger operates on double-entry bookkeeping: every debit
has a matching credit.  Color charge is a ledger quantum number;
a single quark (color-triplet entry) has no matching anti-color
credit.  When two color-carrying sources are separated by distance
$r$, the ledger stores the field energy of the color flux tube:
$E = \sigma r$ (see Section~\ref{sec:potential} for the RS
derivation of $\sigma$).  As $r \to \infty$, this energy diverges:
$J_\mathrm{free} = \lim_{r\to\infty}\sigma r = +\infty$.
Only configurations with zero net color charge (singlets) have
finite $J$-cost and are therefore observable.

\emph{Note}: this argument is structural; it uses the Cornell string tension
$\sigma$ (derived in Section~\ref{sec:potential}) and the RS
identification of color charge with ledger debt.  The result is
self-consistent: the ledger-balance principle motivates confinement,
which is then expressed quantitatively through the Cornell potential.
\end{proof}

\begin{corollary}[Allowed Hadron States]
Color singlets include:
\begin{enumerate}[label=(\roman*)]
  \item \textbf{Mesons}: $q\bar{q}$
  ($3 \otimes \bar{3} \supset 1$).
  \item \textbf{Baryons}: $qqq$
  ($3 \otimes 3 \otimes 3 \supset 1$, via $\epsilon_{ijk}$).
  \item \textbf{Glueballs}: Color-neutral gluon composites.
  \item \textbf{Exotic hadrons}: Tetraquarks ($q\bar{q}q\bar{q}$),
  pentaquarks ($qqqq\bar{q}$), etc., provided they are color
  singlets.
\end{enumerate}
Free quarks ($3$ alone) and free gluons ($8$ alone) are forbidden.
\end{corollary}

%% ============================================================
\section{The Cornell Potential}
\label{sec:potential}
%% ============================================================

\begin{definition}[Confining Potential]
\begin{equation}
  V(r) = -\frac{\alpha_s}{r} + \sigma\,r,
\end{equation}
where $\alpha_s = 2/17$ (see companion paper on $\alpha_s$) and
$\sigma$ is the string tension.
\end{definition}

\begin{proposition}[Confinement at Long Distance]
\label{thm:confining}
For large $r$, $V(r) \approx \sigma r \to \infty$, preventing
quark separation.
\end{proposition}

\begin{proof}
For $r \gg \alpha_s/\sigma$, the Coulomb term $-\alpha_s/r$ is
negligible: $V(r) \approx \sigma r$.  Energy grows without bound.
At $r \sim 1$~fm, $\sigma r \approx 280$~MeV $\approx 2m_\pi$,
triggering $q\bar{q}$ pair creation (string breaking), so isolated
quarks are never observed.
\end{proof}

\begin{proposition}[String Tension --- Structural Estimate]
\label{thm:sigma}
In RS-native units (energy unit $E_{\mathrm{coh}} = \varphi^{-5}$,
length unit $\ell_0$), the string tension has the structural form:
\begin{equation}
  \sigma_{\mathrm{RS}} = \frac{E_{\mathrm{coh}}}{\ell_0}
  = \frac{\varphi^{-5}}{\ell_0}.
\end{equation}
With the hadronic calibration $E_{\mathrm{coh}} \approx 90\;\mathrm{MeV}$
and $\ell_0 \approx 0.2\;\mathrm{fm}$ (matching the QCD scale
$\Lambda_{\mathrm{QCD}}^{-1} \approx 1\;\mathrm{fm}$), the SI value is
$\sigma_{\mathrm{RS}} \approx 0.90\;\mathrm{GeV/fm} \approx 0.18\;\mathrm{GeV}^2$.
Note: the SI value requires the external calibration anchor; see
Remark~\ref{rem:sigma-status}.
\end{proposition}

\begin{remark}[Status of the string tension estimate]
\label{rem:sigma-status}
The identification $\sigma_{\mathrm{RS}} = E_{\mathrm{coh}}/\ell_0$
is a \emph{structural estimate}: the string tension in RS is the
coherence quantum per unit length, the minimum energy stored per
unit of quark separation in the discrete ledger.  The numerical
agreement $\sqrt{\sigma_{\mathrm{RS}}} \approx 420\;\mathrm{MeV}$
vs.\ the lattice QCD value $\sqrt{\sigma} = 440 \pm 10\;\mathrm{MeV}$
\cite{Bali2001} (a 5\% discrepancy) is encouraging.  However, the
precise mapping from RS-native units to SI requires fixing the RS
calibration ($E_{\mathrm{coh}}$ and $\ell_0$); the 5\% agreement
should be understood as order-of-magnitude confirmation rather than
a precision prediction.  A rigorous first-principles derivation of
$\sigma$ from the RS color flux-tube dynamics is an identified open problem.
\end{remark}

%% ============================================================
\section{Mass Gap}
\label{sec:gap}
%% ============================================================

\begin{proposition}[Mass Gap --- Lower Bound from Discreteness]
\label{thm:gap}
The glueball spectrum has a positive lower bound:
$m_G \geq E_{\mathrm{coh}} = \varphi^{-5} > 0$.
\end{proposition}

\begin{proof}[Structural argument]
The discrete ledger has a minimal excitation quantum
$\Delta E \geq E_{\mathrm{coh}} = \varphi^{-5} > 0$.  Any
color-singlet composite above the vacuum (including the lightest
$J^{PC} = 0^{++}$ glueball) requires at least one $\varphi$-rung,
so $m_G \geq E_{\mathrm{coh}} > 0$.  The gap survives the
continuum limit because $\ell_0$ is a physical scale, not a
UV regulator.
\end{proof}

\begin{remark}[Status of the RS mass gap bound]
\label{rem:gap-status}
Lattice QCD predicts the lightest $0^{++}$ glueball at
$m_G \approx 1.7\;\mathrm{GeV}$~\cite{Morningstar1999}.
The RS lower bound ($m_G \geq E_{\mathrm{coh}} \approx 0.09\;\mathrm{GeV}$)
is satisfied by a large margin ($\sim 20 \times$), which reflects
that the RS ledger discreteness argument provides the existence of a
gap, not its quantitative value.  A sharp RS prediction of $m_G$ would
require computing the glueball mass from the RS color singlet structure
and $\varphi$-ladder dynamics---an identified open problem.

This paper does not claim to resolve the Yang--Mills mass gap Millennium
Prize problem; the existence of a gap in RS follows from ledger
discreteness, but this is not the same as proving a mass gap for
four-dimensional Yang--Mills theory in the rigorous mathematical sense
required by the Clay Institute.
\end{remark}

%% ============================================================
\section{Predictions}
\label{sec:predictions}
%% ============================================================

\begin{enumerate}
  \item \textbf{No free quarks}: Quark searches should always
  yield null results.

  \item \textbf{$\sqrt{\sigma} \approx 420\;\mathrm{MeV}$}:
  Consistent with lattice QCD.

  \item \textbf{Mass gap $> 0$}: The glueball spectrum is gapped.

  \item \textbf{$N_c = 3$}: No deviation from three colors.
\end{enumerate}

%% ============================================================
\section{Conclusion}
\label{sec:conclusion}
%% ============================================================

Quark confinement in RS is ledger balance: only color-singlet
(balanced) entries have finite $J$-cost.  Free quarks are
infinitely costly.  The confining potential has Cornell form with
$\sigma = E_{\mathrm{coh}}/\ell_0$, and the mass gap follows from the
discreteness of the ledger.

%% ============================================================
\begin{thebibliography}{99}

\bibitem{Aczel1966}
J.~Acz\'el, \emph{Lectures on Functional Equations and Their Applications},
Academic Press, New York (1966).

\bibitem{RS_Axioms}
J.~Washburn, ``The Algebra of Reality: A Recognition Science Derivation
of Physical Law,'' \emph{Axioms} \textbf{15}(2), 90 (2026).
\texttt{doi:10.3390/axioms15020090}

\bibitem{Wilson1974}
K.~G.~Wilson, ``Confinement of quarks,'' Phys.\ Rev.\ D
\textbf{10}, 2445 (1974).

\bibitem{Bali2001}
G.~S.~Bali, ``QCD forces and heavy quark bound states,''
Phys.\ Rep.\ \textbf{343}, 1 (2001).

\bibitem{Morningstar1999}
C.~J.~Morningstar and M.~J.~Peardon, ``Glueball spectrum from
an anisotropic lattice study,'' Phys.\ Rev.\ D \textbf{60},
034509 (1999).

\end{thebibliography}

\end{document}
